\begin{tikzpicture}[
  x=1mm,y=1mm,
  >=Stealth,
  font=\figurefont\fontsize{8}{10}\selectfont
]

\tikzset{
  axis/.style={->,draw=FigureStroke,line width=.65pt},
  tick/.style={draw=FigureStroke,line width=.6pt},
  ptblue/.style={circle,draw=FigureBlue,fill=FigureBlueFill,minimum size=3.2mm,inner sep=0pt,line width=.6pt},
  ptgreen/.style={circle,draw=FigureGreen,fill=FigureGreenFill,minimum size=3.2mm,inner sep=0pt,line width=.6pt},
  ptyellow/.style={circle,draw=FigureYellow,fill=FigureYellowFill,minimum size=3.2mm,inner sep=0pt,line width=.6pt},
  ptred/.style={circle,draw=FigureRed,fill=FigureRedFill,minimum size=3.2mm,inner sep=0pt,line width=.6pt},
  note/.style={font=\FigureNoteFont\fontsize{7.4}{9}\selectfont,text=FigureInk},
  muted/.style={font=\FigureNoteFont\fontsize{6.9}{8.4}\selectfont,text=FigureMuted}
}

% Number line illustrating the Cheeger-NCut sandwich chain:
% \mu_2 <= OPT_N <= NCut(S_N) <= NCut(S_\phi) <= 2\phi(S_\phi) <= 2\sqrt{2\mu_2} <= 2\sqrt{2 OPT_N}
\draw[axis] (0,8) -- (98,8);

% Positions on axis
\node[ptblue] (p1) at (8,8) {};
\node[ptgreen] (p2) at (23,8) {};
\node[ptgreen] (p3) at (38,8) {};
\node[ptyellow] (p4) at (54,8) {};
\node[ptyellow] (p5) at (69,8) {};
\node[ptred] (p6) at (86,8) {};

\node[note,text=FigureBlue,above] at (8,10.5) {$\mu_2(\mathcal L)$};
\node[muted,below] at (8,5) {谱松弛下界};

\node[note,text=FigureGreen,above] at (23,10.5) {$\operatorname{OPT}_N$};
\node[muted,below] at (23,5) {离散最优};

\node[note,text=FigureGreen,above] at (38,10.5) {$\operatorname{NCut}(S_N)$};
\node[muted,below] at (38,5) {NCut扫描解};

\node[note,text=FigureYellow,above] at (54,10.5) {$2\phi(S_\phi)$};
\node[muted,below] at (54,5) {$\geq\operatorname{NCut}(S_\phi)$};

\node[note,text=FigureYellow,above] at (69,10.5) {$2\sqrt{2\mu_2}$};
\node[muted,below] at (69,5) {Cheeger上界};

\node[note,text=FigureRed,above] at (86,10.5) {$2\sqrt{2\operatorname{OPT}_N}$};
\node[muted,below] at (86,5) {最坏近似上界};

% Braces / arrows showing the two key bridges
\draw[<->,draw=FigureGreen,line width=.55pt] (8,17) to[out=20,in=160] node[above,muted,text=FigureGreen] {谱松弛 $\mu_2\leq\operatorname{OPT}_N$} (23,17);
\draw[<->,draw=FigureRed,line width=.55pt] (38,17) to[out=18,in=162] node[above,muted,text=FigureRed] {Cheeger不等式保证 $\leq 2\sqrt{2\mu_2}$} (69,17);
\end{tikzpicture}
